% !TEX TS-program = lualatex
\documentclass[tikz, border=2mm]{standalone}

% Hỗ trợ Font Tiếng Việt cho LuaLaTeX
\usepackage{fontspec}
\IfFontExistsTF{Times New Roman}{
  \setmainfont{Times New Roman}
}{
  \setmainfont{Liberation Serif}
}
\IfFontExistsTF{Courier New}{
  \setmonofont{Courier New}[Scale=MatchLowercase]
}{}

\usetikzlibrary{shapes.geometric, arrows.meta, positioning, fit, backgrounds, calc}

\begin{document}

% Sơ đồ 1: Giao diện PB2 (Cốt truyện Tương tác)
\begin{tikzpicture}[
    box/.style={draw, rounded corners=3pt, minimum width=5.2cm, minimum height=0.9cm, align=center, font=\small},
    label/.style={font=\footnotesize\itshape, text=gray}
]
    \draw[thick, rounded corners=6pt] (-0.2,-0.3) rectangle (5.4,7.5);
    \node[label] at (2.6, 7.15) {Màn hình 9:16 (dọc)};
    \node[box, fill=blue!8, minimum height=2.6cm] (logo) at (2.6,5.5)
        {\begin{tabular}{c}
            \textbf{Biểu trưng}\\
            \footnotesize Lưới 3$\times$3 ô màu\\
            \footnotesize + chữ C T E T R I S
         \end{tabular}};
    \node[box, fill=green!12, minimum height=0.6cm] (bar) at (2.6,3.9)
        {\textbf{Thanh tải dữ liệu} \ {\footnotesize (Xanh lá, 180$\times$12 px)}};
    \node[label] at (2.6, 3.2) {$\longleftarrow$ khoảng trống nền tối $\longrightarrow$};
    \node[box, fill=gray!10, minimum height=0.7cm] (credit) at (2.6,2.0)
        {\textbf{Dòng ghi công} \ {\footnotesize (Đơn vị phát triển)}};
    \draw[->, gray, dashed] (5.0,5.5) -- (5.8,5.5) node[right, label] {hiện dần};
    \draw[->, gray, dashed] (5.0,3.9) -- (5.8,3.9) node[right, label] {điền dần};
    \draw[->, gray, dashed] (5.0,2.0) -- (5.8,2.0) node[right, label] {hiện dần};
\end{tikzpicture}

% Sơ đồ 2: Hai chế độ vận hành (Story)
\begin{tikzpicture}[
    node distance=0.6cm and 1.4cm,
    block/.style={draw, rounded corners=4pt, minimum width=4.2cm, minimum height=1.5cm, align=center, font=\small},
    cond/.style={draw, diamond, aspect=2.5, align=center, font=\small, fill=yellow!20},
    arr/.style={->, >=Stealth, thick}
]
    \node[cond] (q) {Lần đầu mở ứng dụng?};
    \node[block, fill=blue!8,  below left=of q]  (A)
        {\textbf{Mở đầu + Đồng bộ}\\[3pt]
         \footnotesize Biểu trưng 8 giây\\
         \footnotesize + kiểm tra nội dung mới};
    \node[block, fill=orange!10, below right=of q] (B)
        {\textbf{Chỉ Hội thoại}\\[3pt]
         \footnotesize Bỏ qua biểu trưng\\
         \footnotesize Chạy hội thoại cốt truyện};
    \node[draw, rounded corners=3pt, below=1.8cm of q, minimum width=3cm, fill=green!10, font=\small] (C)
        {Phân hệ Bảng điều khiển};
    \draw[arr] (q) -- node[above left,  font=\footnotesize]{Có} (A);
    \draw[arr] (q) -- node[above right, font=\footnotesize]{Không (chọn truyện)} (B);
    \draw[arr] (A) -- (C); \draw[arr] (B) -- (C);
\end{tikzpicture}

% Sơ đồ 3: Hệ thống Hội thoại
\begin{tikzpicture}[
    box/.style={draw, rounded corners=3pt, minimum width=5.4cm, align=center, font=\small}
]
    \draw[thick, rounded corners=6pt] (-0.2,-0.2) rectangle (5.6,7.6);
    \node[box, fill=gray!15, minimum height=3.4cm] (img) at (2.7,5.5)
        {\textit{Ảnh bối cảnh}\\{\footnotesize (Hiển thị hình ảnh từ máy chủ)}};
    \node[box, fill=blue!10, minimum height=0.5cm] (spk) at (2.7,3.5)
        {\footnotesize\textbf{Tên nhân vật}};
    \node[box, fill=white, minimum height=1.4cm] (txt) at (2.7,2.5)
        {\footnotesize Dòng thoại hiển thị ở đây.\\
         \footnotesize Bấm Xác Nhận để tiếp tục.};
    \node[box, fill=green!8, minimum height=0.5cm, minimum width=2.4cm, dashed] at (1.5,1.1)
        {\footnotesize Lựa chọn A};
    \node[box, fill=orange!8, minimum height=0.5cm, minimum width=2.4cm, dashed] at (3.9,1.1)
        {\footnotesize Lựa chọn B};
    \node[font=\footnotesize\itshape, text=gray] at (2.7, 0.45)
        {(chỉ xuất hiện khi đoạn thoại có rẽ nhánh)};
\end{tikzpicture}

% Sơ đồ 4: Kiểm tra điều kiện sau 8 giây
\begin{tikzpicture}[
    node distance=0.5cm and 1.2cm,
    block/.style={draw, rounded corners=4pt, minimum width=3.6cm, minimum height=1.1cm, align=center, font=\small},
    cond/.style={draw, diamond, aspect=2.2, align=center, font=\small, fill=yellow!20},
    arr/.style={->, >=Stealth, thick}
]
    \node[cond] (c1) at (0,0) {Truyện mới\\mở khoá?};
    \node[block, fill=orange!10] (popup) at (-3,-2.5)
        {\textbf{Hộp thoại Đồng Ý / Từ Chối}\\ \footnotesize ``Chơi truyện mới?''};
    \node[block, fill=blue!8] (console) at (3,-2.5)
        {\textbf{Phân hệ Bảng điều khiển}\\ \footnotesize (Trình đơn chính)};
    \node[block, fill=green!10] (story) at (-5,-4.8)
        {\textbf{Chạy truyện mới}\\ \footnotesize (Chỉ hội thoại)};
    \node[block, fill=gray!10] (replay) at (-1,-4.8)
        {\textbf{Chơi lại truyện cũ}\\ \footnotesize (Chỉ hội thoại)};
    \draw[arr] (c1) -- node[above left,  font=\footnotesize]{Có} (popup);
    \draw[arr] (c1) -- node[above right, font=\footnotesize]{Không} (console);
    \draw[arr] (popup) -- node[left, font=\footnotesize]{Đồng ý} (story);
    \draw[arr] (popup) -- node[right, font=\footnotesize]{Từ chối}  (replay);
\end{tikzpicture}

% Sơ đồ 5: Phân hệ Bảng điều khiển (Trung tâm Điều phối)
\begin{tikzpicture}[
    mod/.style={draw, rounded corners=5pt, minimum width=3cm, minimum height=0.9cm, align=center, font=\small},
    arr/.style={->, >=Stealth, thick}, node distance=0.3cm and 1.5cm
]
    \node[mod, fill=gray!12]  (gs) {Phân hệ Cốt truyện\\{\footnotesize (mở đầu)}};
    \node[mod, fill=blue!10, right=1.5cm of gs] (gc) {\textbf{Bảng điều khiển}\\{\footnotesize (trung tâm)}};
    \node[mod, fill=green!10, right=1.5cm of gc] (co) {Lõi trò chơi\\{\footnotesize (ván đấu)}};

    \draw[arr] (gs) -- node[above,font=\footnotesize]{Hoàn tất} (gc);
    \draw[arr] (gc) -- node[above,font=\footnotesize]{Bấm CHƠI} (co);
    \draw[arr] (co.south) to[bend left=30] node[below,font=\footnotesize]{Thoát ván} (gc.south);
\end{tikzpicture}

% Sơ đồ 6: Giao diện Phiên bản 2 (Bảng điều khiển)
\begin{tikzpicture}[
    zone/.style={draw, rounded corners=3pt, minimum width=5.4cm, align=center, font=\small},
    btn/.style={draw, rounded corners=2pt, minimum width=4.8cm, minimum height=0.52cm, align=center, font=\footnotesize}
]
    \draw[thick, rounded corners=6pt] (-0.2,-0.3) rectangle (5.6,8.2);
    \node[zone, fill=gray!15, minimum height=0.8cm] at (2.7,7.55) {\textbf{C TETRIS -- GAME CONSOLE}};
    \foreach \lbl/\col/\y in {
        HƯỚNG DẪN/blue!20/6.1,
        BẢNG XẾP HẠNG/green!25/5.45,
        CHƠI/red!20/4.8,
        CỐT TRUYỆN/orange!22/4.15,
        CÀI ĐẶT/violet!18/3.5,
        THOÁT/gray!22/2.85
    }{
        \node[btn, fill=\col] at (2.7,\y) {\textbf{\lbl}};
    }
    \node[zone, fill=gray!8, minimum height=1.0cm] at (2.7,1.2) {{\footnotesize TAB / ENTER / ESC -- gợi ý phím tắt}};
\end{tikzpicture}

% Sơ đồ 7: Giao diện Phiên bản 1 (Lõi trò chơi)
\begin{tikzpicture}[font=\small]
    \draw[thick, rounded corners=4pt] (0,0) rectangle (8.4,5.2);
    \draw[fill=gray!8] (0,0) rectangle (5.8,5.2);
    \node[align=center] at (2.9,2.6)
        {\textbf{Bảng chơi}\\240$\times$480 px\\
         \footnotesize 10 cột $\times$ 20 hàng\\
         \footnotesize khối Tetromino rơi ở đây};
    \draw[fill=blue!6] (5.8,0) rectangle (8.4,5.2);
    \node[align=center, font=\footnotesize] at (7.1,4.7) {\textbf{Thanh Bên}};
    \foreach \lbl/\y in {
        {0 THOÁT}/3.95,
        {1 TẠM DỪNG}/3.55,
        {2 ĐIỂM SỐ}/3.15,
        {3 ĐỒNG HỒ}/2.75,
        {4 KẾT TIẾP-1}/2.35,
        {5 KẾT TIẾP-2}/1.95,
        {6 KẾT TIẾP-3}/1.55,
        {7 LẬT CÙNG}/1.15,
        {8 LẬT NGƯỢC}/0.75,
        {9 TRÁI}/0.35
    }{
        \draw (5.85,\y-0.18) rectangle (8.35,\y+0.18);
        \node[font=\tiny] at (7.1,\y) {\lbl};
    }
    \node[font=\tiny] at (7.1,-0.05) {10 PHẢI / 11 TỐC ĐỘ};
\end{tikzpicture}

% Sơ đồ 8: Giao diện Phiên bản 2 (Lõi trò chơi - Hiệu ứng)
\begin{tikzpicture}[
    sbox/.style={draw, rounded corners=3pt, minimum width=2.8cm, minimum height=0.8cm, align=center, font=\footnotesize, fill=white},
    arr/.style={->, >=Stealth, thick},
    node distance=0.3cm
]
    \node[sbox, fill=gray!12] (a) {Hàng đầy\\phát hiện};
    \node[sbox, fill=yellow!20, right=0.5cm of a] (b) {Chớp\\trắng/tối\\$\times$3};
    \node[sbox, fill=green!15,  right=0.5cm of b] (c) {Xoá thật\\+ cộng điểm};
    \node[sbox, fill=blue!10,   right=0.5cm of c] (d) {Khối phía\\trên rơi xuống};

    \draw[arr] (a) -- (b); \draw[arr] (b) -- (c);
    \draw[arr] (c) -- (d);
    \node[font=\tiny\itshape, below=0.15cm of b] {6 nhịp $\times$ 80 ms $\approx$ 480 ms};
\end{tikzpicture}

% Sơ đồ 9: Giao diện Phiên bản 2 (Khối kế tiếp)
\begin{tikzpicture}[font=\small,
    slot/.style={draw, rounded corners=3pt, minimum width=3.8cm, minimum height=1.1cm, align=center}
]
    \node[slot, fill=green!12] (n1) at (0,0)
        {\textbf{KẾ TIẾP-1} (ô 4)\\{\footnotesize Luôn hiện}};
    \node[slot, fill=orange!12] (n2) at (6,0)
        {\textbf{KẾ TIẾP-2} (ô 5)\\
         {\footnotesize Mở khi: Truyện yêu cầu}\\
         {\footnotesize \textbf{và} \texttt{Điểm $\geq$ Ngưỡng}}};
    \node[slot, fill=red!10] (n3) at (12,0)
        {\textbf{KẾ TIẾP-3} (ô 6)\\
         {\footnotesize Mở khi: KẾ TIẾP-2 đã mở}\\
         {\footnotesize \textbf{và} \texttt{Điểm $\geq$ Ngưỡng$\times$2}}};
    \draw[->, >=Stealth, dashed] (n1.east) -- (n2.west) node[midway, above, font=\footnotesize\itshape]{đạt ngưỡng};
    \draw[->, >=Stealth, dashed] (n2.east) -- (n3.west) node[midway, above, font=\footnotesize\itshape]{đạt ngưỡng};
\end{tikzpicture}

% Sơ đồ 10: Giao diện Phiên bản 2 (Lưới cốt truyện)
\begin{tikzpicture}[font=\small]
    \draw[fill=gray!4] (0,2.2) rectangle (3.2,5.5);
    \node[text=gray, font=\footnotesize\itshape] at (1.6,3.85) {(trống)};
    \draw[fill=blue!15]   (0,1.5) rectangle (0.8,2.2);
    \draw[fill=red!15]    (0.8,1.5) rectangle (1.6,2.2);
    \draw[fill=green!15]  (1.6,1.5) rectangle (2.4,2.2);
    \draw[fill=yellow!20] (2.4,1.5) rectangle (3.2,2.2);
    \draw[fill=green!15]  (0,0.7) rectangle (0.8,1.5);
    \draw[fill=blue!15]   (0.8,0.7) rectangle (1.6,1.5);
    \draw[fill=gray!25]   (1.6,0.7) rectangle (2.4,1.5);
    \draw[fill=red!15]    (2.4,0.7) rectangle (3.2,1.5);
    \draw[thick] (0,0.7) rectangle (3.2,5.5);
    \draw[<-, thick] (3.35, 1.45) -- (4.2, 1.45) node[right, font=\footnotesize] {Hàng điền sẵn từ Truyện};
    \draw[<-, gray]  (3.35, 3.85) -- (4.2, 3.85) node[right, font=\footnotesize\itshape] {Phần trống bình thường};
\end{tikzpicture}

\end{document}